\documentclass[11pt]{article}
\usepackage[spanish]{babel}
\usepackage[utf8]{inputenc}
\usepackage[T1]{fontenc}
\usepackage{geometry}
\usepackage{graphicx}
\usepackage{booktabs}
\usepackage{amsmath}
\usepackage{xcolor}
\usepackage{hyperref}
\geometry{margin=2.5cm}

\newcommand{\maybeincludegraphics}[2][]{%
  \IfFileExists{#2}{\includegraphics[#1]{#2}}{\fbox{\parbox{0.88\linewidth}{Figura no disponible: \texttt{#2}}}}%
}

\title{Comunidades topológicas e incompletitud observacional en el catálogo de exoplanetas}
\author{Proyecto Mapper/TDA de exoplanetas}
\date{\today}

\begin{document}
\maketitle

\section{Introducción}
Los sesgos por método de descubrimiento son una propiedad conocida de los catálogos de exoplanetas. Aquí se usan las comunidades Mapper para preguntar si ciertas regiones topológicas parecen submuestreadas por la función de selección observacional. El objetivo no es afirmar planetas faltantes confirmados, sino identificar candidatos a incompletitud observacional.

\section{Motivación astronómica}
El método de tránsito favorece periodos cortos y radios detectables, porque requiere alineación geométrica y tránsitos repetidos. La velocidad radial favorece planetas masivos o señales radiales fuertes. La imagen directa favorece planetas grandes, jóvenes, calientes y separados de su estrella. Microlensing corresponde a una ventana geométrica particular y no debe interpretarse como una clase física directa.

\section{De la auditoría de sesgo a la sombra observacional}
La auditoría previa detectaba asociación entre nodo Mapper y método de descubrimiento. Este análisis pregunta dónde esa asociación, al compararse con el vecindario topológico, sugiere una posible frontera de selección y una región topológica submuestreada.

\section{Datos y grafos analizados}
El análisis se centra en \texttt{orbital\_pca2\_cubes10\_overlap0p35}. Cuando existen artefactos compatibles, también resume \texttt{phys\_min\_pca2\_cubes10\_overlap0p35}, \texttt{joint\_no\_density\_pca2\_cubes10\_overlap0p35}, \texttt{joint\_pca2\_cubes10\_overlap0p35} y \texttt{thermal\_pca2\_cubes10\_overlap0p35}.

\section{Metodología}
Para cada nodo Mapper $v$ y método $m$ se define
\[
p_v(m)=\frac{n_v(m)}{\sum_m n_v(m)}.
\]
El vecindario topológico es $N(v)=\{u:(u,v)\in E\}$ y su composición observacional:
\[
p_{N(v)}(m)=\frac{\sum_{u\in N(v)} n_u(m)}{\sum_m\sum_{u\in N(v)}n_u(m)}.
\]
El contraste local se resume con
\[
D_v(m)=p_{N(v)}(m)-p_v(m),\qquad
R_v(m)=\frac{p_v(m)+\epsilon}{p_{N(v)}(m)+\epsilon}.
\]
La frontera composicional se mide con
\[
\mathrm{method\_l1\_boundary}(v)=\sum_m |p_v(m)-p_{N(v)}(m)|.
\]
El índice principal es
\[
\mathrm{Shadow}(v)=f_{\mathrm{top}}(v)\,[1-H_{\mathrm{norm}}(v)]\,[1-I(v)]\,B(v)\,S(v),
\]
donde $f_{\mathrm{top}}$ es la fracción del método dominante, $I$ la fracción media de imputación, $B$ la distancia L1 al vecindario y $S(v)=\log(1+n_v)/\log(1+\max n)$.

\section{Resultados}
\begin{figure}[ht]
\centering
\maybeincludegraphics[width=0.92\linewidth]{../../outputs/observational_shadow/figures/orbital_graph_by_shadow_score.pdf}
\caption{Grafo orbital coloreado por \texttt{shadow\_score}. Valores altos sugieren una posible frontera de selección local.}
\end{figure}

\begin{figure}[ht]
\centering
\maybeincludegraphics[width=0.92\linewidth]{../../outputs/observational_shadow/figures/orbital_graph_by_shadow_class.pdf}
\caption{Clasificación heurística de nodos por sombra observacional.}
\end{figure}

\begin{figure}[ht]
\centering
\maybeincludegraphics[width=0.76\linewidth]{../../outputs/observational_shadow/figures/orbital_shadow_vs_imputation.pdf}
\caption{Relación entre sombra observacional e imputación media.}
\end{figure}

\begin{figure}[ht]
\centering
\maybeincludegraphics[width=0.76\linewidth]{../../outputs/observational_shadow/figures/orbital_method_boundary_vs_shadow.pdf}
\caption{Sombra observacional frente al contraste local de métodos.}
\end{figure}

\begin{figure}[ht]
\centering
\maybeincludegraphics[width=0.76\linewidth]{../../outputs/observational_shadow/figures/orbital_physical_neighbor_distance_vs_shadow.pdf}
\caption{Sombra observacional frente a distancia físico-orbital al vecindario.}
\end{figure}

\begin{figure}[ht]
\centering
\maybeincludegraphics[width=0.82\linewidth]{../../outputs/observational_shadow/figures/config_shadow_comparison.pdf}
\caption{Comparación resumida entre configuraciones disponibles.}
\end{figure}

\section{Comunidades candidatas a incompletitud}
La tabla \texttt{top\_shadow\_candidates.csv} lista las comunidades con mayor \texttt{shadow\_score}. Cada fila debe leerse como candidato a incompletitud observacional, no como prueba de una población no observada.

\begin{figure}[ht]
\centering
\maybeincludegraphics[width=0.86\linewidth]{../../outputs/observational_shadow/figures/orbital_top_shadow_candidates.pdf}
\caption{Comunidades principales por índice de sombra observacional.}
\end{figure}

\section{Discusión}
Sí se puede afirmar que algunos nodos combinan alta pureza por método, baja entropía, contraste fuerte con sus vecinos y baja dependencia de imputación. No se puede afirmar que esas comunidades contienen planetas no detectados ni que representan una clase planetaria real.

\section{Limitaciones}
Mapper no prueba planetas faltantes. Un vecino topológico no equivale a un planeta real no detectado. Sin funciones de completitud instrumental no se pueden estimar cantidades absolutas. Nodos pequeños pueden tener pureza artificialmente alta. La imputación puede inflar o distorsionar algunas regiones.

\section{Conclusión}
Mapper no solo permite auditar sesgos del catálogo; también puede usarse para priorizar regiones topológicas donde la muestra observada probablemente está incompleta.

\section{Trabajo futuro}
El siguiente paso riguroso sería incorporar funciones de completitud por método, simulaciones de inyección-recuperación o reconstrucciones contrafactuales que permitan separar detectabilidad, arquitectura orbital y estructura física.

\end{document}
